\section{The Sacraments of Initiation}

\begin{massbox}{One supernatural life: birth, strengthening, food}
Baptism, Confirmation, and the Eucharist are not three religious milestones laid side by side. They are one ordered initiation into the Paschal mystery: the font engrafts a person into the crucified and risen Christ; chrism seals and strengthens that baptismal configuration by the Pentecostal Gift; the altar feeds the living member with Christ himself and draws the whole organism into his self-offering. Their proximate effects differ because one life must be born, brought toward mature operation, and sustained by food. Their final end is one: perfected charity in the communion of the Trinity (CCC 1212; ST III, q.~65, a.~1).
\end{massbox}

\subsection{The organism of the New Law}

Aquinas adopts the analogy between bodily and spiritual life while refusing to make it the whole proof of the number seven. Natural life has generation, growth, nourishment, healing from interior fault, healing from exterior debility, public ordering, and propagation. The sacramental economy answers these conditions by Baptism, Confirmation, Eucharist, Penance, Anointing, Orders, and Matrimony. Yet supernatural life is not simply another biological layer. Its formal principle is grace; its governing virtue is charity; its Head is the risen Christ; and its terminus exceeds every created proportion in the vision of God (ST III, q.~65, a.~1; I--II, q.~62, a.~1).

The three initiatory sacraments therefore possess a genuine causal order. Baptism is the doorway because an unbaptized person is not yet capable of receiving the Church's other mysteries. Confirmation presupposes but does not modify the baptismal character: it adds its own distinct character and perfects the baptized for public confession and ecclesial combat. The Eucharist does not imprint another deputation because it gives the end toward which deputation exists: Christ substantially present, offered and received, forming his members into one ecclesial body. In Aquinas's phrase, nearly all the other sacraments are completed in the Eucharist (ST III, q.~65, a.~3).

\begin{threestable}{Initiatory act}{Created effect}{Christological and ecclesial direction}
Birth from water and Spirit & A new creature receives sanctifying grace, faith, hope, charity, and the baptismal character. & Death and resurrection with Christ; incorporation into the Church; capacity for Christian worship.\\
Sealing with chrism & The baptized is strengthened by the Holy Spirit and marked for a more perfect public share in Christ's mission. & Pentecost extends through the apostolic Church; the baptized is distinctly strengthened for public witness.\\
Eating the one Bread & The communicant receives Christ himself; charity and ecclesial unity are increased. & The baptized and sealed body is joined to the Head's sacrifice and given a pledge of resurrection.\\
\end{threestable}

\subsection{Creation, Exodus, Pasch, and Pentecost}

The sensible signs gather the history of salvation rather than abandon it. Water is creation's fertile deep, the flood's judgment and new beginning, Israel's passage through the sea, water from the rock, Jordan, and the opened side of Christ. Oil is royal, priestly, and prophetic consecration; its fragrance signifies the Spirit's invisible diffusion and the ecclesial body anointed in the Anointed One. Bread and wine are creation received through human work, Melchizedek's offering, manna and covenant sacrifice, wisdom's banquet, Passover fulfilled, and the wedding supper still to come. The signs are not arbitrary illustrations subsequently attached to grace. Christ assumes histories, substances, bodily gestures, and words into actions fitted to what he causes.

Patristic mystagogy ordinarily moves across these layers without confusing them. The \emph{Mystagogical Catecheses} traditionally attributed to Cyril of Jerusalem take neophytes from the font through chrism to the Eucharistic table; Ambrose reads creation, Exodus, Naaman, Christ's Baptism, and the Passion within the baptismal water; Augustine returns from the one Eucharistic bread to the ecclesial unity of many grains. Their typology depends upon a realism of divine providence: the same Wisdom who orders creation and covenant brings their signs to fulfillment in the incarnate Word and his Body.

\subsection{The Paschal and Trinitarian grammar}

Every sacrament is an action of the glorified Christ in the Holy Spirit toward the Father. The Father is first origin and final end; the Son's humanity is the conjoined instrument through which the Passion possesses saving efficacy; the Spirit interiorly vivifies what the visible rite outwardly signifies. This order is neither three separate causal streams nor an impersonal mechanism. The one divine operation reaches the creature according to the personal missions: the Son is sent in flesh, the Spirit into hearts, and the adopted creature returns through the Son in the Spirit to the Father (Gal. 4:4--7; ST I, q.~43; III, q.~62, a.~5).

Aquinas consequently calls a sacrament memorial, demonstrative, and prognostic. It remembers the Passion as meritorious cause; demonstrates grace presently effected; and foretells future glory. Initiation makes this threefold time unusually visible: Baptism sacramentally buries and raises; Confirmation communicates the Spirit of the risen Lord; the Eucharist is the memorial sacrifice in which the coming banquet is already pledged. Christian existence is thus stretched between a definitive Pasch and an awaited manifestation: already configured, still being perfected; already fed with immortality, still groaning for bodily redemption.

\begin{studybox}{A teleological guardrail}
The Eucharist is the sacramental end of initiation, but sacramental reception is not the absolute last end. Christ present under the species orders the pilgrim Church beyond sacramental veils to face-to-face vision, when sign yields to possession and God is ``all in all'' (1~Cor. 13:12; 15:28). The sacraments belong to the Church's journey; the charity they cause remains in the homeland.
\end{studybox}
